\section*{\S~3 Abstimmungen und Wahlen}
\label{sec:abstimmungen_und_wahlen}
\addcontentsline{toc}{section}{\S~3 Abstimmungen und Wahlen}

\begin{enumerate}[label=(\arabic*)]
    \item Abstimmungen und Wahlen müssen Allgemein, Unmittelbar, Frei, Geheim und Gleich sein.
    \item Abstimmungen und Wahlen finden entweder im Rahmen der Mitgliederversammlung statt, oder werden von der Leitung, vom Stab, oder vom S1-Feldwebel angesetzt.
    \item Für eine Abstimmung oder Wahl muss ein Wahlleiter bestimmt sein.
    \begin{enumerate}[label=(\Alph*)]
        \item\label{itm:wahlleitung_fall1} Eine Person, die zur Wahl steht, kann nicht die Funktion des Wahlleiters übernehmen. Wenn der Posten des S1-Feldwebels besetzt ist, übernimmt dieser, wenn nicht von mindestens ein Drittel der anwesenden Mitglieder anderes gefordert, bei allen Abstimmungen und Wahlen die Funktion des Wahlleiters. Ist der Posten des S1-Feldwebels nicht besetzt, übernimmt die Unit-Leitung, also Unit-Leiter und Spieß zusammen, die Wahlleitung.
        \item\label{itm:wahlleitung_fall2} Steht eine Wahl zur Unit-Leitung an und der Posten des S1-Feldwebels ist nicht besetzt, oder die Person, die den Posten des S1-Feldwebels innehat, steht zur Wahl, so wird durch die Mitgliederversammlung die Funktion des Wahlleiters bestimmt.
        \item Die Mitgliederversammlung kann auf Verlangen von mindestens ein Drittel der anwesenden Mitglieder verlangen, dass ein anderer Wahlleiter als in \S~3 Absatz 3\ref{itm:wahlleitung_fall1} und 3\ref{itm:wahlleitung_fall2} bestimmt wird. Dies geschieht durch eine öffentliche Abstimmung.
    \end{enumerate}
    \item Abstimmungen und Wahlen müssen mindestens 72 Stunden vor Abstimmungsschluss angekündigt werden, wenn diese nicht unter die Mitgliederversammlung fallen.
    \item Abstimmungen und Wahlen können über eine Drittplattform (bspw. strawpoll.com), Abstimmungen im Discord, einem Hammelsprung in Arma 3, oder durch ein anderes geeignetes Verfahren durchgeführt werden.
    \item Damit eine Abstimmung oder Wahl außerhalb der Mitgliederversammlung beschlussfähig ist, müssen sich mindestens ein Drittel aller Mitglieder beteiligen. Die Beschlussfähigkeit muss nicht überprüft werden.
    \item Abstimmungen können nur die Möglichkeiten „Ja“, „Nein“, oder „Enthaltung“ (oder andere Formulierungen dieser Möglichkeiten) besitzen. Enthaltungen werden bei der Berechnung des Ergebnisses berücksichtigt.
    \item Wahlen müssen immer personenbezogen sein.
    \item Wahlen und Abstimmungen, die keiner anderen Regelung unterliegen brauchen zum bestehen eine relative Mehrheit.
    \item Kann keine Entscheidung getroffen werden, findet eine auf 72 Stunden beschränkte Wahl über einen öffentlichen Kanal der 1st Aircat Division statt. Die Kandidaten für diese Wahl müssen sich direkt im Anschluss an die gescheiterte Wahl bei der Wahlleitung melden. Kann auch hier keine Entscheidung getroffen werden, so wird nach diesem System so lange gewählt, bis es eine Entscheidung gibt.
    \item Bei einer außerordentlichen Wahl übernimmt die gewählte Person den Posten mit Bekanntgabe des Ergebnisses. Die Amtszeit endet zum regulären Zeitpunkt, sofern keine Sonderreglungen existieren.
\end{enumerate}